\author{OTS}
\documentclass[fontsize=16pt]{mybib}
\renewcommand{\DFA}{12mm}
\renewcommand{\DZA}{1.3}

\graphicspath{{../../../assets/images/}}
\setheaderlogo{mnr-green.png}

% Absatzformatierung
\setlength{\parindent}{0pt}
\setlength{\parskip}{0.6em}
\renewcommand{\arraystretch}{0.8}
\newtheorem{theorem}{Theorem}

\title{Glaubwürdigkeit der Schrift}
\author{Lothar Schmid}
\date{}

\begin{document}
\mytitlepage
\tableofcontents
\newpage

\section{Einleitung}
\begin{block}[Folie 1; Begrüssung]
Herzlich willkommen zu dem Vortrag: «Ist die Bibel ein Märchenbuch?»
Ich bin Lothar und einer der Leiter dieser Gemeinde vom Mitternachtsruf hier in Naters. Wer und was der Mitternachtsruf ist, darauf werde ich noch kommen.

Ich muss euch warnen, natürlich bin ich nicht ganz neutral gegenüber der Bibel. Ich glaube, dass die Bibel das Wort Gottes ist. Das ist jetzt nicht einfach so vom Himmel gefallen, sondern durch lesen und erleben, wurde ich ich überzeugt, dass das nicht ein gewöhnliches Buch ist.\\
Aber ich denke, niemand ist wirklich neutral gegenüber der Bibel und vor allem dem Inhalt in diesem Buch.

Zuerst mal zu den Themen die wir zusammen anschauen wollen. 
\begin{enumerate}
    \item Einleitung und Begrüssung
    \item Einzigartigkeit der Bibel
    \item Wie wurde die Bibel hergestellt
    \item Der Kanon
    \item Fazit
\end{enumerate}

Ich muss ehrlich sagen, das ganze Thema ist unglaublich riesig. Als ich mit dem recherchieren anfing, wurde ich richtig überhäuft von den Informationen. Auch waren, Janina und ich im Januar in Berlin an einem vier tägigen Seminar, das dieses Thema zum Schwerpunkt hat. Ich kann in diesem Vortrag nur leicht an der Oberfläche kratzen.\\
Mit jedem einzelne Punkt den wir gemeinsam Anschauen könnte man ein ganzer Tag füllen. Und das ich noch lange nicht das Ende.\\
Der nächste Schritt wäre zu erfahren, ob der Inhalt überhaupt stimmt, gibt widersprüche, gibt es Fehler in der Bibel. Aber dafür reicht leider ein Abend nicht. Wenn erwünscht, kann ich gerne noch einen weiteren Abend planen.
Aber fangen nun an:
\end{block}
\begin{block}[Folie 2; Mitternachtsruf]
    Zuerst einmal will ich den Mitternachtsruf vorstellen:
    \begin{itemize}
        \item Ist eine evangelische Missionsgesellschaft.\\
        Der Mitternachtsruf ist ein evangelisches Missionswerk mit Sitz in Dübendorf Zürich und hat verschiedene Missionsstationen vor allem in Südamerika, Amerika aber auch in Europa.
        \item Logo symbolisiert ein Radiosender\\
        Das Logo symbolisiert einen Radiosender mit einem Radiosender habe sie in Südamerika angefangen. Es wurde um Mitternacht gesendet, weil die Sendezeit dort am billigsten war.        
        \item Der Name kommt aus Matthäus 25:6 \\
        Der Name selber kommt aus Matthäus 25:6 „Um Mitternacht entstand ein Geschrei: Siehe, der Bräutigam kommt!“ Dieses Zitat stammt aus dem Gleichnis mit den 7 törichten und den 7 klugen Jungfrauen und der Bräutigam ist Jesus. Dieser Ruf soll andeuten, dass Jesus eines Tages wiederkommt.
        \item Hauptsitz ist in Dübendorf
        \item Gemeinde in Bern
        \item Gemeinde in Naters\\
        und jetzt seit 3 Jahren hier in Naters. Wie das zustande kam, ist eine lange Geschichte, leider zu lang für heute Abend.
    \end{itemize}
    Also nichts Mystisches nur eine christliche Organisation, die Jesus Christus als Retter verkündigt.
\end{block}
\begin{block}[Folie 3; Zeitschriften]
    Mitternachtsruf vertreibt zwei Gratiszeitschriften.\\
    \textbf{Mitternachtsruf:} Kommentare zum Zeitgeschehen, Prophetie und Auslegungen der Bibel\\
    \textbf{Israel:} Neben den Aktuelle Nachrichten direkt aus Israel selber, wird in dieser Zeitschrift auch die Geschichte, die Kultur und ihre Feste erklärt.

    Die Zeitschriften könnt ihr hinten beim Eingang gratis Abonnieren, indem ihr die Karte ausfüllt über den QR-Code auf die Webseite geht oder euch einfach ein Probeexemplar mitnehmen.

So jetzt geht’s los
\end{block}
\begin{block}[Folie 4; Übersicht der Bibel]
    \textbf{Wie beginnt das Märchenbuch}\\
    Wie beginnt ein Märchenbuch?
Oft beginnt es mit den Worten \betonung[b, p]{„Es war einmal…“} Früher irgend wann mal, irgendwo, hat dieses Märchen stattgefunden oder nicht. Jeder kann sich selber eine Zeit aussuchen und sich dies zurecht legen. Und irgendwo hat es stattgefunden, in der Schweiz, in Deutschland vielleicht weil der Autor von Hönsel und Gretel aus Deutschland kommt? Wir können uns das selber definieren und zurchtlegen. Gewöhnlich liegt ein Moral dahinter, wo jetzt bei Rotköppchen ist, kann ich nicht sagen. Bei Hönsel und Gretel darin, dass man nicht alleine in den Wald gehen soll.
\textbf{Wie beginnt die Bibel}\\
Im Anfang schuf Gott die Himmel und die Erde.
Es ist im Anfang. Es gibt kein Vorher. Die Aussage ist klar. 
So fängt die Bibel an. Die Bibel ist ein Buch in dem Geschichte und Moral vereint ist. Tausende haben schon versucht, Fehler zu finden, Geschichte zu widerlegen, bis jetzt konnte noch niemand etwas vorbringen was gegen die Geschichte spricht. Auch die Wunder in diesen Buch, wir können sagen, Heute sind wir aufgeklärt und wissen es besser, aber noch konnte kein Mensch etwas erschaffen, was nicht schon vorher da war. Wir wollen das nur nicht gerne zugeben.

Aber schauen wir mal was die Bibel so faszinierend macht.
\end{block}
\begin{block}[Folie 5; Allgemeines]
    Bilbios ist ein griechiches Wort und heisst eigentlich nur Buch oder Schriftrolle. Die antike griechische Hafenstadt Byblos war ein Umschlagplatz für Papyrus und daher stammt wohl der Name.
Im lateinischen wurde der Name dann zu biblia und im Deutschen zur Bibel. Also übersetzt einfach nur Buch. Die Bibel ist eine Sammlung von Bücher und Schriften.

Das Buch ist geteilt ins AT mit 39 Büchern und ins NT mit 27 Büchern. 
Es beginnt mit 1. Mose Kapitel 1 und endet mit der Offenbarung Kapitel 22
Von der Erschaffung bis zur Auflösung der Erde und Erschaffung einer neuen Erde.

Ein paar Fakten dazu:

\end{block}
\begin{block}[Folie 6, 7; Fakten]
    \begin{enumerate}
        \item Geschrieben über eine Zeitspanne von 1400 Jahren\\
        ca. 1300 \vChr{} bis 90\nChr{}
        \item Geschrieben von mehr als 40 Autoren aus verschiedenen Gesellschaftschichten
        \item Geschrieben an verschiedenen Orten:
        \item Geschrieben zu verschiedenen Zeitspannen
        \item Geschrieben in verschiedenen Gemütsverfassungen\\
        Wenn wir Paulus nehmen der bei seinen letzten Briefen angekettet an einen römischen Soldaten in Rom im Gefängnis sass.\\
        Oder Moses der 40 Jahre mit einem störrischen Volk durch die Wüste wanderte. Können wir uns vorstellen, dass jeder mit anderen Emotionen schrieb.
        \item Geschrieben in drei Sprachen
        Das Alte Testament auf Hebräsich und teile in Aramäisch, eine dem hebräischen verwandte Sprache und das Neue Testament in Griechisch.
    \end{enumerate}
\end{block}
\begin{block}[Folie 8; Prophetie]
    Die Bibel ist das einzige bekannte Buch mit nachweisbar erfüllten Prophetien oder Voraussagen. 
Vor allem über das Volk Israel, Jesus seine Geburt und sein Sterben und über unsere Zukunft.

Hier ein Beispiel aus dem Buch Jesaja Kapitel 45 Vers 1.
\begin{bibelbox}{SCHL}{Jes}{45:1}
So spricht der Herr zu Kyrus, seinem Gesalbten, dessen Rechte Hand ich ergriffen habe, um Völker vor ihm niederzuwerfen und die Lenden der Könige zu entgürten, um Türen vor ihm zu öffnen und Tore, damit sie nicht geschlossen bleiben:
\end{bibelbox}
Der Kyrus der hier erwähnt wird, wurde erst 150 Jahre nachdem dies hier aufgeschrieben geboren. \betonung[b, p]{seinem Gesalbten}, ein Gesalbter ist ein König. Hier sagt Gott, dass er den persischen König Kyrus für seine Zwecke auserkoren hat. Kyrus wird dann auch der König sein, der den Juden erlaubt nach der Gefangenschaft in Babylon zurück nach Israel zu gehen und den Tempel in Jerusalem wieder aufzubauen. Wie schon gesagt, dass wurde schon 150 Jahre vor der Geburt des Königs mit Namen aufgeschrieben.


Jesaja 53 wird von den Rabbis in der Synagoge nicht vorgelesen, weil dort konkrete Angaben über Jesu und sein Sterben gemacht werden. 
\begin{bibelbox}{SCHL}{Jes}{53:3}
Verachtet war er und verlassen von den Menschen, ein Mann der Schmerzen und mit Leiden vertraut; wie einer, vor dem man das Angesicht verbirgt, so verachtet war er, und wir achteten ihn nicht.
\end{bibelbox}
Lange Zeit wurde gesagt, dass dieser Text erst nach Jesus geschrieben wurde, bis man 1946 in Qumran am Toten Meer zwei vollständige Jesaja Rollen fand, die man auf 200 \vChr{} datierte.

Noch ein Beispiel aus den Psalmen. Die Psalmen sind eine Sammlung von Gebeten und Liedern, vor allem von König David, der den Goliath erschlagen hat. Dieser Schreibt:
\begin{bibelbox}{SCHL}{Ps}{22:19}
Sie teilen meine Kleider unter sich und werfen das Los über mein Gewand.
\end{bibelbox}
Und im Neuen Testament steht 1000 Jahre später:
\begin{bibelbox}{SCHL}{Matt}{27:35}
Nachdem sie ihn nun gekreuzigt hatten, teilten sie seine Kleider unter sich und warfen das Los\dots
\end{bibelbox}

Wir könnten noch lange so fortfahren. Roger Liebi ein Schweizer Theologe, Hystoriker studierte unter anderem die alten Sprachen, sagt, dass es im Alten Testament rund 300 erfüllte Prophetien oder voraussagen über Jesus gibt.
\end{block}
\begin{block}[Folie 9; Sie hat die Zeit überlebt]
Die Bibel ist ja aktuell um die 3300 Jahre alt. Sie hat die Zeit überlebt, obwohl sie fast 3000 Jahre von Hand kopiert werden musste und dafür vergängliches Material wie Papyrus und Pergament benutzt wurde.
Es gibt immer noch tausende von Manuskripten, der grösste Fund war in Qumran am Toten Meer. Heute 80 Jahre später sind sie immer noch dabei, die gefundenen Fragmente zu übersetzen und zu archivieren. 

Als wir in Israel waren, haben wir einen Fotografen kennengelernt und er hat uns gezeigt wie sie heute mit den modernsten Methoden das Alter und das Material analysieren.
Die Funde sind überigends nicht geheim wie manche sagen, sondern das archivierte Material kann man auf der Webseite anschauen und wer Lust hat kann beim Übersetzen mithelfen.
\end{block}
\begin{block}[Folie 11; Sie hat Verfolgung überlebt]
\end{block}
\end{document}